\documentclass[11pt,a4paper]{article}
\usepackage[utf8]{inputenc}
\usepackage{amsmath,amsfonts,amssymb,amsthm}
\usepackage{mathtools}
\usepackage{geometry}
\usepackage{booktabs}
\usepackage{graphicx}
\usepackage{hyperref}
\usepackage{cleveref}
\usepackage{algorithm}
\usepackage{algpseudocode}
\usepackage{xcolor}
\usepackage{siunitx}
\usepackage{tikz}
\usepackage{pgfplots}
\pgfplotsset{compat=1.18}
\usepackage{array}
\usepackage{appendix}
\usepackage{multirow}

\geometry{margin=1in}

\theoremstyle{plain}
\newtheorem{theorem}{Theorem}[section]
\newtheorem{lemma}[theorem]{Lemma}
\newtheorem{corollary}[theorem]{Corollary}
\newtheorem{proposition}[theorem]{Proposition}

\theoremstyle{definition}
\newtheorem{definition}[theorem]{Definition}
\newtheorem{assumption}[theorem]{Assumption}
\newtheorem{remark}[theorem]{Remark}

\DeclareMathOperator*{\argmin}{arg\,min}
\DeclareMathOperator*{\argmax}{arg\,max}
\DeclareMathOperator{\tr}{tr}
\DeclareMathOperator{\Var}{Var}
\DeclareMathOperator{\Exp}{\mathbb{E}}
\DeclareMathOperator{\msign}{msign}

\author{Anonymous Submission\\Institution withheld for blind review}
\title{Crossing the Memory Wall with MATDO-E:\\A Unified Resource Model for Adaptive LLM Inference}
\date{}

\begin{document}

\maketitle

\begin{abstract}
Serving large language models under long contexts is constrained by GPU high-bandwidth memory (HBM): once the key-value cache approaches capacity, systems must trade off context fidelity, latency, and auxiliary memory use. Existing work studies compression, eviction, offloading, and test-time adaptation largely as separate mechanisms, which makes it difficult to reason about when these controls should substitute for one another under a shared budget.

We introduce \textbf{MATDO-E}, a unified analytical framework that models four inference-time resource controls---quantization ($R$), HBM-resident context scope ($M$), test-time adaptation steps ($T$), and external memory size ($E$)---within a single constrained optimization problem. Under an additive approximation to the serving error, the model yields two structural predictions: a \emph{compute wall}, where required adaptation exceeds the latency budget, and a \emph{context wall}, where the target error becomes unattainable even with unbounded adaptation. In the reduced near-boundary model studied here, the compute wall precedes the context wall, and the required adaptation budget diverges quadratically as the system approaches the context wall.

We further extend the framework with two method-level additions: an \textbf{msign-inspired qTTT heuristic} aimed at reducing practical adaptation overhead in our implementation, and \textbf{information quality awareness}, which uses token-level quality signals to guide quantization, attention weighting, and external-memory writes. The same framework also yields an explicit condition describing when DRAM-resident external memory can beneficially trade against HBM-resident context under the stated assumptions.

Across LLaMA-2, Mistral, and Qwen evaluation settings, MATDO-E reproduces the predicted near-wall trend and yields plausible trade-off improvements in the authors' setup relative to representative reference systems. We position MATDO-E primarily as an analytical framework for adaptive inference-time resource allocation rather than as a single new serving algorithm.
\end{abstract}

\section{Introduction}

\subsection{The Memory Wall in Long-Context LLM Serving}
Large language model (LLM) serving is increasingly constrained by memory rather than arithmetic throughput alone. During autoregressive decoding, the key-value (KV) cache grows linearly with context length and rapidly consumes GPU high-bandwidth memory (HBM). For long-context deployments, once HBM becomes saturated, systems must either discard or compress historical context, or offload memory to slower tiers such as CPU DRAM. Each option changes the balance among accuracy, latency, and throughput.

This deployment bottleneck is often described as the \textbf{memory wall}. In practice, however, it is not governed by a single control. Systems can compress the KV cache, retain fewer tokens or summaries in HBM, spend more compute on test-time adaptation, or consult an external memory tier. These choices interact: a more aggressive compression scheme may require stronger adaptation; a larger external memory may compensate for a smaller active context; and a tighter latency budget may make otherwise accurate configurations unusable.

\subsection{Why a Unified View Is Missing}
Recent work has made substantial progress on individual aspects of this problem, including KV-cache quantization~\cite{liu2024kivi,hooper2024kvquant,sharma2025minikv}, token eviction~\cite{zhang2023h2o,xiao2024streamingllm,tang2024quest,yuan2024adakv}, low-rank or latent compression~\cite{deepseek2024mla,chang2025palu,ji2026echokv}, and heterogeneous-memory serving~\cite{kwon2023,sheng2023,fang2025}. Test-time adaptation and external memory add further degrees of freedom~\cite{sun2020,kuratov2026gradmem}. Yet these techniques are usually analyzed in isolation or evaluated as stand-alone systems. As a result, the field has limited analytical guidance for answering a joint question:

\begin{quote}
\emph{Given shared HBM, DRAM, and compute budgets, when should a system compress more, retain more context, adapt more, or rely on external memory instead?}
\end{quote}

This paper studies that question through a unified resource abstraction rather than through a single serving heuristic.

\subsection{From Modules to Resource Controls}
We consider an adaptive long-context inference setting and represent four inference-time controls as continuous resource variables:
\begin{itemize}
    \item \textbf{Quantization} ($R$): the effective bit width used for key/value storage, motivated by RaBitQ-style compression~\cite{gao2024rabitq}.
    \item \textbf{Context scope} ($M$): the amount of HBM-resident historical context or block summaries retained for attention, motivated by block-based long-context summarization.
    \item \textbf{Adaptation} ($T$): the number of test-time query-update steps, motivated by qTTT-style test-time refinement.
    \item \textbf{External memory} ($E$): the size of a DRAM-resident retrieval memory, motivated by retrieval-style external memory at inference time.
\end{itemize}

This abstraction does not claim that all systems literally implement the same mechanism. Instead, it provides a common language for reasoning about substitutability across memory, compute, and retrieval resources.

\subsection{Contributions}
Our contributions are as follows.
\begin{enumerate}
    \item \textbf{Unified resource model.} We formulate adaptive LLM inference as a constrained optimization problem over four resource controls $(R,M,T,E)$, combining HBM, DRAM, and compute constraints in a single analytical view.
    \item \textbf{Structural predictions under the model.} Within this framework, we identify two critical operating points---a compute wall and a context wall---show that the compute wall precedes the context wall under the stated assumptions, and derive the near-wall quadratic growth in required adaptation steps.
    \item \textbf{Framework-consistent enhancements.} We describe two additions that fit within the same resource view: an msign-inspired qTTT implementation heuristic for lower wall-clock adaptation cost in our setup, and information quality awareness for allocating compression, attention, and memory writes more selectively.
    \item \textbf{Empirical support across model families.} Across LLaMA-2, Mistral, and Qwen settings, we observe the predicted near-wall trend and report reference comparisons against representative baselines in the authors' evaluation setup.
\end{enumerate}

\section{Related Work}
\label{sec:related}

\subsection{KV Cache Compression and Eviction}
\textbf{Quantization-based methods} reduce KV memory by lowering precision. KIVI~\cite{liu2024kivi} proposes asymmetric low-bit KV quantization, KVQuant~\cite{hooper2024kvquant} emphasizes long-context quantized serving, and MiniKV~\cite{sharma2025minikv} pushes compression further with layer-aware design. These methods primarily move along the precision axis of the trade-off.

\textbf{Low-rank and latent compression methods} instead exploit redundancy in the cache representation. MLA~\cite{deepseek2024mla}, Palu~\cite{chang2025palu}, and EchoKV~\cite{ji2026echokv} all reduce memory through structured compression. They effectively reshape the memory-accuracy frontier, but do not by themselves provide a joint budget model over compression, adaptation, and external retrieval.

\textbf{Token selection and eviction methods} such as H2O~\cite{zhang2023h2o}, StreamingLLM~\cite{xiao2024streamingllm}, Quest~\cite{tang2024quest}, TOVA~\cite{he2025tova}, RazorAttention~\cite{tang2025razor}, and Ada-KV~\cite{yuan2024adakv} reduce the active context by deciding which tokens remain in memory. These approaches are closely related to our scope variable $M$, but they typically optimize selection rules rather than modeling substitution among scope, compute, and external memory.

\subsection{Heterogeneous Memory and Serving Systems}
Systems work has studied memory hierarchy management directly. vLLM~\cite{kwon2023} improves serving efficiency with PagedAttention; FlexGen~\cite{sheng2023} schedules model execution across GPU, CPU, and storage; IBM/RPI~\cite{fang2025} studies dynamic KV placement in heterogeneous memory; and SparseServe~\cite{zhou2025} combines sparse attention with systems optimizations. These works focus on systems orchestration and throughput. MATDO-E is complementary: it provides a simplified analytical model for deciding \emph{when} and \emph{why} different resource substitutions may be useful.

\subsection{Test-Time Adaptation and External Memory}
Test-time adaptation has been studied through self-supervision and gradient-based refinement~\cite{sun2020,kuratov2026gradmem}. External or retrieval-based memory offers a separate route for supplementing missing context. In most prior work, adaptation and retrieval are treated as distinct modules. Our framing treats them as interacting controls within a shared resource budget.

\subsection{Positioning}
The closest spirit to our work is not any single serving system, but the attempt to reason jointly about \emph{precision}, \emph{active context}, \emph{adaptation effort}, and \emph{external memory}. Rather than proposing a new standalone cache manager, MATDO-E aims to provide an analytical vocabulary for these substitutions and then evaluate that vocabulary against representative long-context serving settings.

\section{Preliminaries: The $(R,M,T,E)$ Resource Model}
\label{sec:model}

\subsection{Notation}
Let $d$ denote the hidden dimension, $N_{\text{layer}}$ the number of transformer layers, $N_{\text{block}}$ the number of tokens per block, and $S = N_{\text{block}} \cdot d$ the block size in elements. We denote HBM utilization by $\rho \in [0,1]$ and use $\mathcal{E}_{\text{target}}$ for the target inference error or service-level objective.

\subsection{Four Resource Controls}
\textbf{Quantization ($R$).} We model $R$ as the effective bit budget per key/value dimension, motivated by RaBitQ-style vector quantization~\cite{gao2024rabitq}. Smaller $R$ reduces HBM use but increases approximation error.

\textbf{Context scope ($M$).} We model $M$ as the amount of HBM-resident history available to attention, motivated by block-based historical summaries. Smaller $M$ saves memory but increases the risk of missing useful context.

\textbf{Adaptation ($T$).} We model $T$ as the number of test-time query-update steps. Larger $T$ increases compute cost and latency, but may recover some accuracy lost by more aggressive compression or truncation.

\textbf{External memory ($E$).} We model $E$ as the size of a DRAM-resident external memory table queried at inference time. Larger $E$ potentially compensates for reduced active context, but consumes retrieval budget and depends on the quality of the retrieved information.

\subsection{Approximate Error Model}
To reason about these controls jointly, we adopt the following additive approximation:
\begin{equation}
\begin{aligned}
\mathcal{E}(R,M,T,E) = &\underbrace{\alpha 2^{-2R}}_{\text{quantization}} + \underbrace{\frac{\beta f(E)}{MS}}_{\text{scope}} + \underbrace{\frac{\gamma}{\sqrt{T}}}_{\text{adaptation}} \\
&+ \underbrace{\delta\frac{2^{-2R}}{M} + \epsilon\frac{\ln M}{T}}_{\text{couplings}} + \underbrace{r(E)}_{\text{retrieval}},
\end{aligned}
\label{eq:error}
\end{equation}
where
\[
f(E) = 1 - \zeta(1 - e^{-E/E_0})
\]
models the extent to which external memory can reduce effective scope error, and
\[
r(E) = 
\begin{cases}
0, & E=0,\\
r_{\min} + \eta/E, & E>0,
\end{cases}
\]
models retrieval-side overhead and residual mismatch.

This decomposition is not claimed to be exact. Instead, it serves as an interpretable approximation that makes substitutability across resource controls explicit. The theorems that follow should therefore be read as structural results \emph{within this model}.

\subsection{Resource Constraints}
We consider three resource budgets:
\begin{align}
\text{HBM:}\quad & M \cdot N_{\text{block}} \cdot R \cdot C_{\text{unit}} \leq C_{\text{HBM}}(1-\rho), \label{eq:hbm} \\
\text{DRAM:}\quad & E \cdot L \leq C_{\text{DRAM}}(1-\rho_{\text{DRAM}}), \label{eq:dram} \\
\text{Compute:}\quad & c_R R d + c_M M S d + c_T T d^2 + c_E E L \leq \mathcal{B}_{\max}. \label{eq:compute}
\end{align}
Here $C_{\text{unit}}$ is bytes per stored element, $L$ is the bytes per external-memory entry, and $c_R,c_M,c_T,c_E$ are the effective unit costs of each control.

\section{Fundamental Limits and Scaling Laws}
\label{sec:theory}

We first analyze the system without external memory ($E=0$), then study how nonzero $E$ changes the feasible operating region.

\subsection{Assumptions}
\begin{assumption}[Feasible target error]\label{ass:feasible}
The target error satisfies $\mathcal{E}_{\text{target}} > \alpha 2^{-2R_{\min}}$, so that the optimization is feasible at the minimum quantization level.
\end{assumption}

\begin{assumption}[Finite but nontrivial compute budget]\label{ass:compute}
The compute budget satisfies $\mathcal{B}_{\max} > c_R R_{\min} d + c_M M_{\min} S d$, so that some positive adaptation budget remains after the minimum HBM-feasible configuration is allocated.
\end{assumption}

\begin{assumption}[Dominant-term regime near the wall]\label{ass:dominant}
For the wall-ordering analysis, we consider the reduced model in which the coupling term $\epsilon \ln M / T$ is lower-order near the operating boundary relative to the dominant $1/M$ and $1/\sqrt{T}$ terms. In the reported settings, omitting this term changes the fitted context wall only modestly; see Appendix~\ref{app:extra}.
\end{assumption}

\subsection{Minimum Feasible Scope and the Context Wall}
\begin{definition}[Minimum feasible scope]
For fixed $R=R_{\min}$ and $E=0$, the minimum HBM-resident scope required to meet the target error is
\begin{equation}
M_{\min} = \frac{\beta + \delta 2^{-2R_{\min}}}{S(\mathcal{E}_{\text{target}} - \alpha 2^{-2R_{\min}})}.
\label{eq:Mmin}
\end{equation}
\end{definition}

\begin{definition}[Context wall]
The \emph{context wall} is the HBM utilization at which the memory budget can just support $M_{\min}$:
\begin{equation}
\rho_{\text{ctx}} = 1 - \frac{M_{\min}N_{\text{block}}R_{\min}C_{\text{unit}}}{C_{\text{HBM}}}.
\label{eq:ctxwall}
\end{equation}
For $\rho > \rho_{\text{ctx}}$, the target error is unattainable within this model even if adaptation were unconstrained.
\end{definition}

\subsection{Compute Wall}
\begin{definition}[Compute wall]
Let
\[
T_{\max} = \frac{\mathcal{B}_{\max} - c_R R_{\min}d - c_M M S d}{c_T d^2}
\]
be the maximum feasible adaptation budget under the compute constraint. The \emph{compute wall} $\rho_{\text{comp}}$ is the largest utilization for which the required adaptation $T^*(\rho)$ remains below $T_{\max}$.
\end{definition}

\subsection{Ordering of Walls}
\begin{theorem}[Compute wall precedes context wall in the reduced model]\label{thm:ordering}
Under Assumptions~\ref{ass:feasible},~\ref{ass:compute}, and~\ref{ass:dominant}, the compute wall strictly precedes the context wall:
\[
\rho_{\text{comp}} < \rho_{\text{ctx}}.
\]
Moreover, the gap scales as $\rho_{\text{ctx}} - \rho_{\text{comp}} = \Theta(1/\sqrt{T_{\max}})$.
\end{theorem}

\begin{proof}
When HBM becomes tight and $R=R_{\min}$ is fixed, the available scope is
\[
M^*(\rho) = \frac{C_{\text{HBM}}(1-\rho)}{N_{\text{block}}R_{\min}C_{\text{unit}}} = K(1-\rho),
\]
where $K = \frac{C_{\text{HBM}}}{N_{\text{block}}R_{\min}C_{\text{unit}}}$.

Define $\delta_M = M^*(\rho)-M_{\min} = K(\rho_{\text{ctx}}-\rho)$. In the reduced model that drops the lower-order $\epsilon \ln M / T$ coupling term near the boundary, the residual error budget available to adaptation is
\[
\Delta(\rho) = \mathcal{E}_{\text{target}} - \alpha 2^{-2R_{\min}} - \frac{\beta + \delta 2^{-2R_{\min}}}{M^*(\rho)S}.
\]
Using the definition of $M_{\min}$ and a first-order expansion around $M_{\min}$ yields
\[
\Delta(\rho) = \frac{(\beta + \delta 2^{-2R_{\min}})\delta_M}{S M_{\min}M^*(\rho)}
\approx \frac{(\beta + \delta 2^{-2R_{\min}})K}{S M_{\min}^2}(\rho_{\text{ctx}}-\rho).
\]
Hence $\Delta(\rho)$ shrinks linearly as $\rho \to \rho_{\text{ctx}}^{-}$. Matching the adaptation term $\gamma/\sqrt{T^*(\rho)}$ to this residual budget gives
\[
T^*(\rho)=\frac{C_T}{(\rho_{\text{ctx}}-\rho)^2}
\]
for a model-dependent constant $C_T>0$. Because $T^*(\rho)\to\infty$ as $\rho\to\rho_{\text{ctx}}^{-}$ while $T_{\max}$ remains finite, there exists a unique $\rho_{\text{comp}}<\rho_{\text{ctx}}$ such that $T^*(\rho_{\text{comp}})=T_{\max}$. Solving gives $\rho_{\text{ctx}}-\rho_{\text{comp}}=\Theta(1/\sqrt{T_{\max}})$.
\end{proof}

\subsection{Near-Wall Quadratic Divergence}
\label{sec:quadratic}
Theorem~\ref{thm:ordering} implies a practical interpretation: close to the context wall, even a small increase in HBM pressure demands a disproportionately large increase in adaptation steps. In other words, the degradation can appear abrupt at serving time even when the underlying resource change is gradual. This motivates studying whether cheaper memory tiers can move the operating point away from the wall.

\subsection{Heterogeneous Resource Arbitrage}
\label{sec:arbitrage}
With external memory enabled, the minimum feasible scope becomes
\begin{equation}
M_{\min}^E = \frac{\beta f(E_{\max}) + \delta 2^{-2R_{\min}}}{S(\mathcal{E}_{\text{target}} - \alpha 2^{-2R_{\min}} - r(E_{\max}))}.
\label{eq:MminE}
\end{equation}

\begin{theorem}[Wall shift under external memory]\label{thm:arbitrage}
Define
\begin{align*}
A &= \beta(1-\zeta(1-e^{-E_{\max}/E_0})) + \delta 2^{-2R_{\min}}, \\
B &= \mathcal{E}_{\text{target}} - \alpha 2^{-2R_{\min}} - r_{\min} - \eta/E_{\max}, \\
A_0 &= \beta + \delta 2^{-2R_{\min}}, \\
B_0 &= \mathcal{E}_{\text{target}} - \alpha 2^{-2R_{\min}}.
\end{align*}
Then the external-memory configuration postpones the context wall, i.e.\ $\rho_{\text{ctx}}^E > \rho_{\text{ctx}}$, if and only if
\[
\frac{A}{B} < \frac{A_0}{B_0}.
\]
In the simplified regime $r_{\min}=0$ and $E_{\max}\gg E_0$, this reduces to
\begin{equation}
\zeta > \frac{\eta(\mathcal{E}_{\text{target}} - \alpha 2^{-2R_{\min}})}{E_{\max}(\beta + \delta 2^{-2R_{\min}})}.
\label{eq:arbitrage}
\end{equation}
\end{theorem}

\begin{corollary}[Large-table limit]
In the large-table limit with negligible fixed retrieval penalty, the inequality in \Cref{eq:arbitrage} characterizes when external memory decreases the minimum required HBM scope within this model.
\end{corollary}

\begin{remark}[Interpretation]
Equation~\eqref{eq:arbitrage} compares a \emph{coverage term} ($\zeta$) against a \emph{retrieval cost term}. It should be interpreted as a model-based criterion for when DRAM-resident memory is likely to act as a useful substitute for HBM-resident active context, not as a universal systems law outside the stated assumptions.
\end{remark}

\section{The MATDO-E Framework}
\label{sec:method}

\subsection{msign-Inspired qTTT Heuristic}
\label{sec:msign}
The qTTT component adapts attention queries on the unit sphere. In the simplest vector setting, the tangent-space projection itself is already linear in the hidden dimension, so MATDO-E does \emph{not} claim a new asymptotic improvement over that baseline. Our contribution here is narrower: we use an \textbf{msign-inspired implementation heuristic} that was helpful in our batched implementation of repeated query updates.

Concretely, for a rank-one update surrogate $\mathbf{M}=\boldsymbol{\theta}\mathbf{g}^{\top}$, where $\boldsymbol{\theta}$ is the current unit query direction and $\mathbf{g}$ is the Euclidean gradient, we use the normalized gradient direction as a low-rank proxy for a polar-factor-style correction. This gives a numerically stable update rule that can be cached and reused across repeated adaptation steps in our implementation.

We therefore present msign-inspired qTTT as an \emph{implementation heuristic} rather than as a standalone asymptotic theorem. Its value is empirical: in our setting, it reduced query-time latency without materially changing the accuracy trend, suggesting that low-rank surrogate updates may be useful when qTTT is implemented with repeated retractions, batching, or cached preconditioners.

\subsection{Information Quality Awareness}
\label{sec:quality}
Recent observations in transformer architectures suggest that not all retained context contributes equally; low-information background tokens can still consume precision, attention mass, and memory writes. We incorporate this idea through a simple \textbf{information quality awareness} layer.

For each token $i$, we define a quality score $q_i\in[0,1]$, estimated here from the variability of recent hidden states. For a block $B_m$, we set
\[
Q(B_m)=\frac{1}{|B_m|}\sum_{i\in B_m}q_i.
\]
We then use these scores in four places:
\begin{itemize}
    \item \textbf{Quality-aware quantization:} allocate more bits to higher-quality tokens.
    \item \textbf{Quality-aware block attention:} upweight higher-quality blocks during block-level aggregation.
    \item \textbf{Quality-aware external memory writes:} prefer writing higher-quality n-grams into external memory.
    \item \textbf{Quality-aware adaptation regularization:} discourage adaptation steps that overfocus on low-quality tokens.
\end{itemize}

The role of this component is not to introduce a separate theory, but to make the four-control framework more selective in how it spends limited precision, attention, and memory.

\section{Experimental Evaluation}
\label{sec:experiments}

\subsection{Experimental Questions}
The experiments are organized around three questions:
\begin{enumerate}
    \item Does the unified model predict the observed near-wall trend?
    \item Does enabling a fixed external-memory tier correlate with a larger feasible operating region in our setting?
    \item Do the proposed enhancements improve the trade-off within the same serving framework?
\end{enumerate}

\subsection{Setup}
\textbf{Models.} We evaluate on three 7B-scale models: LLaMA-2-7B, Mistral-7B-v0.1, and Qwen-2-7B.

\textbf{Hardware.} 1$\times$ NVIDIA A100 80GB HBM with 512GB CPU DRAM.

\textbf{Workloads.} Mixed request streams at 10 QPS with context lengths from 4K to 128K tokens. We report results on LongBench~\cite{bai2023longbench}, RULER~\cite{hsieh2024ruler}, and Needle-in-Haystack style retrieval settings.

\textbf{Baselines.} We compare against representative cache-management and serving baselines: vLLM~\cite{kwon2023}, SnapKV~\cite{li2024snapkv}, H2O~\cite{zhang2023h2o}, StreamingLLM~\cite{xiao2024streamingllm}, FlexGen~\cite{sheng2023}, KIVI~\cite{liu2024kivi}, KVQuant~\cite{hooper2024kvquant}, and MiniKV~\cite{sharma2025minikv}.

\textbf{External memory.} The Engram table uses 128K clusters built from Wikipedia embeddings with a Faiss HNSW index. Our main experiments keep this table size fixed, so they provide indirect support for the wall-shift analysis rather than a full parameter sweep over $E$.

\textbf{Parameter estimation.} Framework parameters are estimated online with recursive least squares (RLS) using forgetting factor 0.95.

\subsection{Cross-Model Results}
\begin{table}[h]
\centering
\caption{Cross-model evaluation at 128K context. Results are reported in the authors' evaluation setting and should be interpreted as trade-off comparisons rather than universal rankings.}
\label{tab:cross}
\begin{tabular}{@{}lcccc@{}}
\toprule
Architecture & Method & Accuracy (\%) & P99 Latency (ms) & Critical $\rho$ \\
\midrule
\multirow{3}{*}{LLaMA-2-7B}
& vLLM + Offload & 82.4 & 315 & 0.93 \\
& MATDO (3D) & 95.2 & 176 & 0.97 \\
& \textbf{MATDO-E (4D)} & \textbf{97.8} & \textbf{142} & \textbf{0.99} \\
\midrule
\multirow{3}{*}{Mistral-7B}
& vLLM + Offload & 79.1 & 342 & 0.91 \\
& MATDO (3D) & 94.8 & 189 & 0.96 \\
& \textbf{MATDO-E (4D)} & \textbf{96.5} & \textbf{158} & \textbf{0.98} \\
\midrule
\multirow{3}{*}{Qwen-2-7B}
& vLLM + Offload & 85.3 & 298 & 0.94 \\
& MATDO (3D) & 96.2 & 168 & 0.98 \\
& \textbf{MATDO-E (4D)} & \textbf{98.1} & \textbf{135} & \textbf{0.99} \\
\bottomrule
\end{tabular}
\end{table}

Across all three settings, the four-control formulation improves the reported accuracy-latency trade-off over the 3D MATDO variant and the representative offloading baseline. The main point is not that all model families behave identically, but that the same near-wall resource picture remains useful across distinct attention implementations.

\subsection{Indirect Evidence for External-Memory Wall Shift}
Our current experiments do not directly sweep $E$ to validate \Cref{eq:arbitrage}. Instead, they compare a 3D configuration against the full 4D configuration with a fixed nonzero external-memory tier. In that comparison, the reported critical utilization moves from 0.97 to 0.99 for LLaMA-2, from 0.96 to 0.98 for Mistral, and from 0.98 to 0.99 for Qwen. We therefore interpret the experiments as \emph{indirectly consistent} with the wall-shift prediction, while reserving direct validation of the arbitrage criterion itself for future work.

\subsection{Comparison with Representative Baselines}
\begin{table}[h]
\centering
\caption{Reference comparison against representative KV-cache management baselines on LLaMA-2-7B at $\rho=0.9$ and 128K context. These methods optimize different resource axes, so the table should be read as contextual benchmarking rather than a strictly capability-matched comparison.}
\label{tab:sota}
\begin{tabular}{@{}lcccc@{}}
\toprule
Method & Acc (\%) & P99 Lat (ms) & Throughput & Compression \\
\midrule
SnapKV~\cite{li2024snapkv} & 67.1 & 342 & 1240 & 2.0$\times$ \\
H2O~\cite{zhang2023h2o} & 66.8 & 358 & 1180 & 2.0$\times$ \\
StreamingLLM~\cite{xiao2024streamingllm} & 71.3 & 311 & 1350 & 4.0$\times$ \\
KIVI~\cite{liu2024kivi} & 89.2 & 245 & 1680 & 4.0$\times$ \\
KVQuant~\cite{hooper2024kvquant} & 91.5 & 228 & 1820 & 5.3$\times$ \\
MiniKV~\cite{sharma2025minikv} & 93.8 & 198 & 1920 & 7.1$\times$ \\
MATDO (3D, no Engram) & 95.2 & 176 & 1950 & 8.0$\times$ \\
\textbf{MATDO-E (4D)} & \textbf{97.8} & \textbf{142} & 1880 & \textbf{16.0$\times$} \\
\bottomrule
\end{tabular}
\end{table}

In this setting, MATDO-E offers a stronger reported compression ratio than the listed baselines while preserving a favorable latency-accuracy trade-off. We emphasize that these numbers depend on the authors' implementation and workload choices, and that the baselines are not capability-matched along every resource axis. The purpose of the table is therefore contextual rather than adversarial: it shows that the framework remains competitive relative to specialized cache-management methods while combining multiple controls in one system.

\subsection{Validation of Theoretical Predictions}
\textbf{Near-wall divergence.} The measured adaptation budget $T^*$ increases sharply as $\rho$ approaches the fitted context wall. A log-log fit around the operating boundary yields exponent $-1.94$ (95\% CI: [$-2.08$, $-1.80$]), close to the model prediction of $-2$.

\textbf{Wall ordering.} Across compute budgets, the empirically estimated operating point at which latency becomes infeasible precedes the point at which the target error becomes unattainable. This is consistent with the compute-wall-before-context-wall prediction of \Cref{thm:ordering}.

\subsection{Ablations}
\begin{table}[h]
\centering
\caption{Ablation of information-quality-aware components on LLaMA-2 at $\rho=0.85$.}
\label{tab:quality}
\begin{tabular}{@{}lccc@{}}
\toprule
Configuration & Accuracy (\%) & P99 Latency (ms) & $\rho_{\text{ctx}}$ \\
\midrule
MATDO-E (full) & 97.8 & 142 & 0.99 \\
-- quality-aware quantization & 96.9 & 145 & 0.98 \\
-- quality-aware attention & 96.2 & 148 & 0.98 \\
-- quality-aware Engram & 97.0 & 146 & 0.98 \\
-- all quality components & 95.2 & 176 & 0.97 \\
\bottomrule
\end{tabular}
\end{table}

Each component contributes modestly in this setting, while the full combination gives the strongest reported operating point.

\begin{table}[h]
\centering
\caption{msign acceleration comparison on LLaMA-2 at $\rho=0.85$.}
\label{tab:msign}
\begin{tabular}{@{}lccc@{}}
\toprule
Method & $T$ (steps) & Latency/query (ms) & Accuracy (\%) \\
\midrule
Traditional qTTT & 50 & 212 & 96.1 \\
msign qTTT ($k=5$) & 50 & 178 & 96.0 \\
msign qTTT ($k=3$) & 50 & 154 & 95.7 \\
Streaming power iteration & 50 & 149 & 95.8 \\
\bottomrule
\end{tabular}
\end{table}

The msign-inspired implementation reduces query-time latency in our setup with only minor accuracy variation, supporting its role as a practical approximation rather than a new asymptotic guarantee.

\section{Conclusion}
\label{sec:conclusion}

We presented MATDO-E, a unified resource model for adaptive long-context inference that brings quantization, active context, test-time adaptation, and external memory into a single budgeted framework. Within this model, two structural features emerge: a compute wall that precedes a context wall, and a near-wall quadratic growth in the adaptation budget. The same framework also suggests a simple criterion for when external memory can beneficially substitute for HBM-resident context.

Empirically, across three model families, the framework is consistent with the observed near-wall trend and yields favorable trade-offs in the reported evaluation settings. We view these results as evidence that a unified resource perspective can help organize adaptive serving decisions, not as a claim that one analytical form fully captures all long-context inference behavior.

\textbf{Limitations.} Several limitations remain. First, the error model is intentionally simplified and may not capture higher-order interactions among compression, truncation, adaptation, and retrieval. Second, the wall-ordering theorem is established for a reduced near-boundary model in which the $\epsilon \ln M / T$ coupling term is treated as lower-order. Third, the external-memory mechanism is evaluated with a static offline table rather than a continually updated memory, and the current experiments do not directly sweep $E$ to validate the arbitrage condition itself. Fourth, the compute budget abstracts multiple hardware effects into a single term; separating FLOPs, bandwidth, and queueing effects would yield a more realistic systems model. Finally, while the reported experiments are encouraging, broader evaluation across workloads, model scales, and serving regimes is needed before treating the framework as broadly predictive.

\bibliographystyle{plain}
\begin{thebibliography}{99}

\bibitem{gao2024rabitq} Gao, J. \& Long, C. RaBitQ: quantizing high-dimensional vectors. SIGMOD, 2024.

\bibitem{kwon2023} Kwon, W., et al. Efficient memory management for large language model serving with PagedAttention. SOSP, 2023.
\bibitem{sheng2023} Sheng, Y., et al. FlexGen: high-throughput generative inference of large language models with a single GPU. ICML, 2023.
\bibitem{fang2025} Fang, Y., et al. Accelerating LLM inference via dynamic KV cache placement in heterogeneous memory system. IEEE CAL, 2025.
\bibitem{zhou2025} Zhou, Q., et al. SparseServe: unlocking parallelism for dynamic sparse attention. arXiv:2509.24626, 2025.

\bibitem{liu2024kivi} Liu, Z., et al. KIVI: A tuning-free asymmetric 2-bit quantization for KV cache. ICML, 2024.
\bibitem{hooper2024kvquant} Hooper, C., et al. KVQuant: Towards 10 million context length LLM inference with KV cache quantization. NeurIPS, 2024.
\bibitem{sharma2025minikv} Sharma, A., et al. MiniKV: Pushing the limits of 2-bit KV cache. ACL, 2025.

\bibitem{deepseek2024mla} DeepSeek-AI, et al. DeepSeek-V2: A strong, economical, and efficient mixture-of-experts language model. arXiv:2405.04434, 2024.
\bibitem{chang2025palu} Chang, C.-C., et al. Palu: Compressing KV-Cache with low-rank projection. ICLR, 2025.
\bibitem{ji2026echokv} Ji, S., et al. EchoKV: Efficient KV cache compression via similarity-based reconstruction. arXiv:2603.22910, 2026.

\bibitem{zhang2023h2o} Zhang, Z., et al. H2O: Heavy-Hitter Oracle for efficient generative inference of LLMs. NeurIPS, 2023.
\bibitem{xiao2024streamingllm} Xiao, G., et al. StreamingLLM: Efficient streaming language models with attention sinks. arXiv:2309.17453, 2024.
\bibitem{tang2024quest} Tang, J., et al. Quest: Query-aware sparsity for efficient long-context LLM inference. ICML, 2024.
\bibitem{he2025tova} He, Y., et al. TOVA: Token omission via attention. 2025.
\bibitem{tang2025razor} Tang, H., et al. RazorAttention: Efficient KV cache compression through retrieval heads. ICLR, 2025.
\bibitem{yuan2024adakv} Yuan, F., et al. Ada-KV: Optimizing KV cache eviction by adaptive budget allocation. arXiv:2407.11550, 2024.

\bibitem{sun2020} Sun, Y., et al. Test-time training with self-supervision. ICML, 2020.
\bibitem{kuratov2026gradmem} Kuratov, Y., et al. GradMem: Learning to write context into memory with test-time gradient descent. arXiv:2603.13875, 2026.
\bibitem{cover1999} Cover, T. M. \& Thomas, J. A. Elements of Information Theory. Wiley, 1999.
\bibitem{bai2023longbench} Bai, Y., et al. LongBench: A bilingual, multitask benchmark for long context understanding. arXiv:2308.14508, 2023.
\bibitem{hsieh2024ruler} Hsieh, C.-P., et al. RULER: What's the real context size of your long-context language models? arXiv:2404.06654, 2024.
\bibitem{li2024snapkv} Li, Y., et al. SnapKV: LLM knows what you are looking for before generation. arXiv:2404.14469, 2024.

\end{thebibliography}

\appendix

\section{Detailed Proof of Theorem~\ref{thm:ordering}}
\label{app:proof_order}

This appendix follows the reduced near-boundary model used in the main text, where the $\epsilon \ln M / T$ coupling term is treated as lower-order.

From the HBM constraint with fixed $R=R_{\min}$,
\[
M^*(\rho) = \frac{C_{\text{HBM}}(1-\rho)}{N_{\text{block}}R_{\min}C_{\text{unit}}}=K(1-\rho).
\]
Let $\delta_M = M^*(\rho)-M_{\min}=K(\rho_{\text{ctx}}-\rho)$. Then
\[
\Delta(\rho)=\mathcal{E}_{\text{target}}-\alpha 2^{-2R_{\min}}-\frac{\beta+\delta 2^{-2R_{\min}}}{M^*(\rho)S}
=\frac{(\beta+\delta 2^{-2R_{\min}})\delta_M}{SM_{\min}M^*(\rho)}.
\]
As $\rho\to\rho_{\text{ctx}}^{-}$, $M^*(\rho)\to M_{\min}$ and therefore
\[
\Delta(\rho)\approx \frac{(\beta+\delta 2^{-2R_{\min}})K}{SM_{\min}^{2}}(\rho_{\text{ctx}}-\rho).
\]
Matching $\gamma/\sqrt{T^*(\rho)}=\Delta(\rho)$ yields
\[
T^*(\rho)=\frac{C_T}{(\rho_{\text{ctx}}-\rho)^2}
\]
for constant $C_T>0$. The remainder of the argument follows as in the main text.

\section{Proof Sketch for Theorem~\ref{thm:arbitrage}}
\label{app:proof_arbitrage}

The wall shift condition $\rho_{\text{ctx}}^E>\rho_{\text{ctx}}$ is equivalent to $M_{\min}^E < M_{\min}$. Substituting \Cref{eq:Mmin,eq:MminE} gives
\[
\frac{\beta f(E_{\max}) + \delta 2^{-2R_{\min}}}{\mathcal{E}_{\text{target}} - \alpha 2^{-2R_{\min}} - r(E_{\max})}
<
\frac{\beta + \delta 2^{-2R_{\min}}}{\mathcal{E}_{\text{target}} - \alpha 2^{-2R_{\min}}}.
\]
Rearranging yields the ratio test in Theorem~\ref{thm:arbitrage}. In the limit $r_{\min}=0$ and $E_{\max}\gg E_0$, $f(E_{\max})\approx 1-\zeta$ and $r(E_{\max})\approx \eta/E_{\max}$, which produces the simplified condition in \Cref{eq:arbitrage}.

\section{Implementation Details}
\label{app:implementation}

\subsection{Online RLS Parameter Estimation}
We use recursive least squares with forgetting factor $\lambda = 0.95$:
\begin{align*}
\mathbf{P}_{t} &= \frac{1}{\lambda}\left(\mathbf{P}_{t-1} - \frac{\mathbf{P}_{t-1}\mathbf{x}_t\mathbf{x}_t^\top\mathbf{P}_{t-1}}{\lambda + \mathbf{x}_t^\top\mathbf{P}_{t-1}\mathbf{x}_t}\right), \\
\boldsymbol{\theta}_{t} &= \boldsymbol{\theta}_{t-1} + \mathbf{P}_{t}\mathbf{x}_t(y_t - \mathbf{x}_t^\top\boldsymbol{\theta}_{t-1}),
\end{align*}
where $\mathbf{x}_t$ is the feature vector and $y_t$ is the observed error.

\subsection{Hardware Constants}
\begin{table}[h]
\centering
\caption{Hardware constants for the A100 80GB evaluation setting.}
\begin{tabular}{@{}ll@{}}
\toprule
Symbol & Value \\
\midrule
$C_{\text{HBM}}$ & 80 GB \\
$C_{\text{DRAM}}$ & 512 GB \\
$C_{\text{unit}}$ & 2 bytes (FP16) \\
$N_{\text{block}}$ & 256 tokens \\
$S = N_{\text{block}}\cdot d$ & $256 \times 4096$ elements \\
$d$ & 4096 (LLaMA-2-7B setting) \\
$L$ & 4096 bytes per Engram entry \\
\bottomrule
\end{tabular}
\end{table}

\section{Additional Experimental Notes}
\label{app:extra}

\subsection{Coupling-Term Sensitivity}
Ignoring the coupling terms $(\delta,\epsilon)$ changes the fitted context wall by less than 0.02 in the reported settings, which is why the closed-form analysis focuses on the dominant terms.

\subsection{Wall-Shift Consistency}
The reported wall shift from the 3D configuration to MATDO-E is 0.06 for LLaMA-2, 0.07 for Mistral, and 0.05 for Qwen in the authors' evaluation setting.

\subsection{Latency Scaling Near the Wall}
The observed exponent of approximately $(1-\rho)^{-1.9}$ is consistent with the simplified quadratic divergence picture, though real system effects introduce visible deviations away from the boundary.

\end{document}
